\section{Coefficient-uniform torus tails and literal three-row primitive germs}
\label{sec:r12-uniform-tail-three-row-germs}

\paragraph{Status and attribution.}
\textbf{PROVED THIS ROUND, submitted for independent exact-source review.}
The constructor is \texttt{/root/r11\_challenge}.  The exact EFT,
EFC, XTF and GNT constructions are frozen inputs.  The new estimate
makes their relevant zero-tail bound independent of the direction
of the original section tuple, including as that tuple approaches
its common zero set.  This permits a finite bounded-parameter
comparison of congruent lifts of the prescribed two-pivot three-row
construction.  No uniform bound is imposed on an entire noncompact
neighborhood, and no arbitrary primitive is given naturality.

\subsection{A coefficient-only tail threshold}

Use the original block rows $2,3$, the external row $1$, a rank
$r\ge3$, and the fixed whole EFC data $(S,v)$ over an ordinary
base open.  Every native cap constant and the canonical GB
partition function are fixed before making comparisons.  The
metric is always the same original metric.  Initially work on
the ordinary overlap of the two Gaussian pivots of $S$.

Fix a base point $x$.  Shrink its ordinary neighborhood $O$ so
that the EFT6 whole coefficient-zero calls have their prescribed
zero expressions there.  The remaining complete GB formula has
the finite norm list EFT10.  Its matrices and inverse matrices
are analytic and bounded after a further local shrinking, and
its positive cap floors have a positive lower bound.  Thus choose
constants
\[
 A\ge\max(1,|b|:b\in\mathcal B),\qquad
 V\ge\max(1,|v_2|,|v_3|),\qquad 0<c\le1
 \tag{UTG1}
\]
on $O$, with $\mathcal B$ the exact finite EFT7 list, such that
\[
 c\le\|F_j^{-1}\|^{-1},\qquad
 c\le\lambda_j\|H_j^{-1}\|^{-1}
                      \quad\hbox{for every retained norm or cap.}
 \tag{UTG2}
\]
No entry of the section tuple occurs in these choices.  The
notation $V$ in (UTG1) denotes a positive real bound, not an
analytic domain or the external shear matrix $V(v)$.

Choose one $T>0$ such that
\[
                         e^{-T}A<1,\qquad e^{-T}2V<c.
 \tag{UTG3}
\]
For an \emph{arbitrary} section tuple $h$ of the given metrized
line bundle over $O$, write
\[
 w=v_2h_2+v_3h_3,\quad R_0=\|(h_2,\ldots,h_r)\|,
 \quad M_1=\max(\|h_1\|,\|h_1+w\|).
 \tag{UTG4}
\]
The following two ordinary opens cover $O\setminus|h=0|$:
\[
 O_h^{(0)}=\{\|w\|<\|h_1\|\},\qquad
 O_h^{(1)}=\{M_1<2VR_0\}.
 \tag{UTG5}
\]
Indeed, outside the first open one has
$\|h_1\|\le\|w\|\le VR_0$, and hence $M_1\le VR_0$.
Here $R_0>0$, unless the entire tuple is zero.  This proves
membership in the second strict open.  If $R_0=0$ and the
tuple is nonzero, the point is in the first open.

On $O_h^{(0)}$ all full norms and capped passive norms for
$D_1(t)h$ and $D_1(t)V(v)h$ agree for every $t$, by the exact
EFT11 argument.  On $O_h^{(1)}$, whenever $\tau=v(t)>T$,
\[
                         |t|M_1<e^{-T}2VR_0<cR_0.
 \tag{UTG6}
\]
The exact positive floors (UTG2) therefore give equality of
every full norm and every capped passive subtuple norm in
the two complete GB expressions.  All other coefficient
weights and partitions are identical.  The same strict tail
also makes both complete unpartitioned crossing triples zero
by EFT6--8, since $|t|A<1$.

\begin{theorem}[A tuple-independent local tail]
On the coefficient neighborhood $O$, the one number $T$ chosen
in (UTG3) works for every section tuple $h$ simultaneously:
\[
 \mathcal T(S,tv;D_1(t)h)=0
 \quad\hbox{on }(O\setminus|h=0|)\times\{\tau>T\}.
 \tag{UTG7}
\]
This is the whole actual triple, including its mixed parameter
components, not only its boundary or its pointwise scalar value.
The same conclusion holds for the entire strict crossing witness:
\[
 q_X(S,tv;D_1(t)h)=0
 \quad\hbox{on the same tail after increasing $T$ if needed.}
 \tag{UTG8}
\]
The constants depend only on the locally fixed coefficient data
and the finite native call list, not on $h$ or a lower bound for
any of its nonzero entries.
\end{theorem}

\begin{proof}
Equations (UTG5)--(UTG6) prove the full norm-and-cap equality
on an ordinary cover, so all actual derivatives agree.  Combining
it with the whole-call zero domains proves (UTG7).
For (UTG8), the inner PTH parameter $s$ in XTF2 has
$v(st)=\sigma+\tau$.  If $\tau>T$ and $\sigma\ge0$, the
whole inner family is zero by (UTG7), for the same arbitrary
tuple $h$.  If $\sigma<0$, the full $K_{\sigma_+}$ operator
is zero.  The input is also zero on a neighborhood of the
common wall $\sigma=0$.  Consequently its actual PTH integral
is zero.  The two unpartitioned terms in EFC4 are already zero
by EFT8.  This is exactly the full XTF proof, now with the
coefficient-only constant (UTG3).
\end{proof}

The assertion is local also across a zero of one Gaussian pivot.
There the canonical GB partition is constant on an ordinary
coefficient neighborhood.  EFC4 reduces to its one surviving
unpartitioned word, while the differentiated overlap term is
identically absent.  Its finite EFT6 list supplies a coefficient-only
bound directly.  No singular overlap matrix is retained there.

For later comparison of the whole TPG face, the same estimate
applies to every finite list of block-multiplication witnesses
from GMU, TML and SGI.  Their norm-and-cap lists are precisely
GNT9.  In (UTG2) take their finite union, retaining the endpoints
of every formal CT transport.  Equality of both endpoint caps
implies equality of the whole interpolated formal cap and every
extracted coefficient, including $-d''J$.  Thus
\[
 q_T(A;D_1(t)V(v)h)=q_T(A;D_1(t)h)
 \quad(\tau>T)
 \tag{UTG9}
\]
with a coefficient-only local $T$.  This statement does not
change the passive block of an ST witness or infer equality
of two triples merely from their boundaries.

\subsection{Uniform positive-floor comparison on bounded parameter slabs}

The equality notion for coefficient lifts is literal equality
of their values in $\mathscr H(x)$ at a chosen original analytic
point $x$.  We compare the same fixed native construction with
the same factor order, row labels, cap offsets and smooth
partition functions.  Corresponding original metrics agree.
If an initial input has the form $F_0f$, the corresponding
initial matrix $\widetilde F_0$ must satisfy
$F_0(x)=\widetilde F_0(x)$ and be invertible near $x$.

Fix a compact real parameter interval $I$.  After the whole
coefficient-zero calls have been removed on their ordinary
zero neighborhoods, every retained parameter matrix is a
finite Laurent matrix in $t$ with base analytic coefficients.
Its determinant and inverse denominators are monomial parameter
units times base units.  This is the exact EFT14 table; the
base matrices in the additional block multiplications have
parameter exponent zero.  It is not a claim for an arbitrary
family containing a unit $1+ta$.

Suppose corresponding base coefficients agree at $x$.
Every retained base unit has exactly the same norm for the
two systems on an ordinary neighborhood of $x$, by
$|\widetilde u-u|<|u|$.  Thus all valuation weights and the
canonical base partitions agree on that whole neighborhood.
Their monomial versions $t^a u$ agree in norm for every $t$.
For each finite matrix pair $F(t),\widetilde F(t)$, equality
at $x$ is equality as Laurent matrices over $\mathscr H(x)$.
The finite upper and lower bounds on $|t|$ for $\tau\in I$
therefore allow one ordinary coefficient neighborhood on which
\[
 \|(\widetilde F(t)-F(t))F(t)^{-1}\|<1
                         \quad\hbox{for every }\tau\in I.
 \tag{UTG10}
\]
For example, expand the finitely many matrix entries and
inverse entries in Laurent powers of $t$.  Each coefficient
difference vanishes at $x$; its product with the finitely
many bounds $\sup_{\tau\in I}|t|^a$ tends to zero near $x$.
This gives (UTG10) simultaneously for the entire finite list.

For a capped subtuple with common floor $\lambda(t)>0$,
the floor has a positive lower bound on a smaller coefficient
neighborhood times $I$: its exponent is a fixed finite
piecewise affine expression in the equal base unit weights
and bounded $\tau$.  Shrink further to impose
\[
 \|\bigl(\widetilde H(t)-H(t)\bigr)F(t)^{-1}\|
                                    <\lambda(t)
                         \quad\hbox{for every }\tau\in I.
 \tag{UTG11}
\]
The original GIC1--6 positive-floor proof now gives, for
\emph{every} tuple $f$ off its common zero set,
\[
 \begin{gathered}
 \|F(t)f\|=\|\widetilde F(t)f\|,\\
 \max\{\|H(t)f\|,\lambda(t)\|F(t)f\|\}
 =\max\{\|\widetilde H(t)f\|,
                       \lambda(t)\|\widetilde F(t)f\|\}.
 \end{gathered}
 \tag{UTG12}
\]
This remains true when any passive subtuple is zero.  No
lower bound for a nonzero component of $f$ was used.

At a coefficient zero, the relevant \emph{whole} zero-call
inequalities remain valid uniformly on the compact slab after
shrinking, using $\sup_{\tau\in I}|t|$ and the bounded retained
base units.  The term is then absent for both data sets.
No singular normalizer is compared separately.  For the
native coefficient-zero extensions of GMU/TML use their
entire local second-pivot expressions, as in GNT9; do not
retain subterms already canceled on that ordinary open.

Every surviving scalar in the prescribed program is now
literally equal on the whole ordinary product neighborhood.
Consequently all actual first derivatives, corner currents,
metric transport terms and mixed jets agree.  Finite formal
CT coefficient extraction preserves the same equality.
This proves bounded-slab equality of the complete actual
families, not just their restrictions to a chosen skeleton.

\subsection{The integrated external witness preserves primitive germs}

Consider two data sets $(S,v,F_0)$ and
$(\widetilde S,\widetilde v,\widetilde F_0)$ with corresponding
coefficient values equal at $x$.  All block determinants and
initial matrices are invertible near $x$.  For the common
original tuple $f$ the following is a literal equality of
punctured germs:
\[
 q_X(S,v;F_0f)
      =q_X(\widetilde S,\widetilde v;\widetilde F_0f).
 \tag{UTG13}
\]
To prove it at a point with two nonzero block pivots, first
apply (UTG7) to both coefficient systems and choose one common
$T$ and neighborhood.  This choice does not depend on $f$.
Use the bounded-slab comparison (UTG10)--(UTG12) with a
slightly larger interval containing $[0,T]$.  On the resulting
ordinary neighborhood the complete two parameter inputs
$\mathcal T^t$ agree on that slab.  They are both actual
zero in the positive tail, whereas $K_{\tau_+}$ is zero
for $\tau<0$.  Hence the full locally slab-supported PTH
integrands agree as actual forms on the entire product.

PTH's entire actual-zero preservation, checked on every base
affinoid unit refinement, makes their actual integrals equal.
The two remaining unpartitioned terms of EFC4 agree by the
same finite comparison.  This proves (UTG13), including the
complete differentiated inner partition term.  At a zero of
one block pivot that term is identically absent and the
surviving native word comparison proves (UTG13) directly.
This is equality of the entire triple, not only a third
boundary or equality of its second boundary.

No common unit presentation through $f=0$ is asserted.  The
comparisons hold on its ordinary punctured complement, where
the forms have their actual local presentations.  Their
neighborhood and their parameter cutoff are independent of
the direction of approach to $f=0$.  This distinction is
essential for comparing conormal lifts.

\subsection{Literal germs of the prescribed two-pivot three-row pair}

Let $M,\widetilde M$ be two invertible three-row blocks as in
TPG2, with $M(x)=\widetilde M(x)$.  Suppose their two specified
first-column entries $a,d$ are not both zero at $x$.  Let
$p_M(F_0f),E_M$ denote exactly the TPG7 construction, with
the same chosen inner and outer partition functions and
the same cap constants for both systems.

\begin{theorem}[Three-row two-pivot germ stability]
On an ordinary neighborhood of $x$,
\[
 p_M(F_0f)=p_{\widetilde M}(\widetilde F_0f)
                    \quad\hbox{off }|f=0|,\qquad
 E_M=E_{\widetilde M}\quad\hbox{on the whole neighborhood}.
 \tag{UTG14}
\]
The statement compares congruent lifts of the same prescribed
program.  It does not compare unrelated Gaussian atlases or
different words for the same matrix.
\end{theorem}

\begin{proof}
If $a(x)=0$ or $d(x)=0$, the outer partition is constant
near $x$ for both systems.  The entire TPG overlap correction
is then absent.  The surviving local pair is a finite SHT
word, one canonical GB subblock, and another finite SHT
word.  Corresponding intermediate matrices agree at $x$.
The exact GIC finite-word germ theorem gives equality of
these pairs and residues.

Suppose now $a(x)d(x)\ne0$.  The two outer valuation ratios,
partitions and all their derivatives agree on an ordinary
neighborhood.  Both local pairs $P_a,P_d$ agree by GIC,
applied to the corresponding Schur blocks and intermediate
matrices.  To compare their full G3F15 overlap triples,
retain its four summands separately.
\begin{enumerate}
\item Every outside call is the same exact ST, DTC, MX, NSC,
      SL, CS or WNR call at congruent invertible intermediate
      matrices.  Nonzero unit weights agree; coefficient-zero
      calls use their whole ordinary zero expressions.
      GIC positive-floor stability gives literal equality of
      every original passive cap and complete MT expression.
\item The strict external crossing summand agrees by (UTG13),
      including its actual parameter integration.
\item The upper-prefix same-target SAN companion has the same
      whole vertices, unit weights and normalization zero.
      Small-coefficient companions are literally zero.
      Its finite formula therefore agrees for both lifts.
\item The finite block-multiplication witness agrees by the
      bounded finite norm-and-cap program GNT9, with every
      endpoint of its nested CT transports retained.  The
      original ordinary zero-pivot extensions in GMU/TML
      are used as whole formulas.
\end{enumerate}
Thus the entire $Q_{ad}$ agrees.  Its extending pair also
agrees on the unpunctured neighborhood: the outside CS error
depends on equal unit valuations, while the remaining GMU/TML
errors have their exact input-independent coefficient formulas
and whole zero-pivot extensions.  At those zero pivots their
second-pivot formulas involve only the retained nonzero units;
the coefficient-zero cancellations are identical for the two
lifts.  No unpunctured equality is inferred by density from
the punctured domain.  The same is true of $E_a,E_d$ by
GIC.  Substitution into the complete TPG7 formula proves
(UTG14) in both domains.
\end{proof}

For a supplied finite word of such two-pivot three-row blocks,
ordinary two-row blocks, diagonal units and shears, the same
result follows term by term at the actual intermediate tuple.
All intermediate matrices for two congruent lifts have equal
values at $x$, and a finite intersection of ordinary neighborhoods
suffices.  The fixed-program pairs are equal as punctured germs;
the extending residues are equal as unpunctured actual germs.

\paragraph{Topology, relations and limits.}
The objects used here are actual local forms and their literal
ordinary germs.  Their local presentations are the supplied finite
native norm/cap expressions, with PTH integrations over a common
bounded parameter interval where needed.  Unit refinements and
subdivisions preserve the scalar equalities before differentiation;
PTH separately preserves all actual zero relations.  On any fixed
coefficient presentation and fixed support slab its integral has
the exact coefficient-final continuity proved in PTH.  No
topological claim about choosing the comparison neighborhood or
varying the whole matrix data is required.  No normed completion,
full-space continuous retraction, or inherited-kernel comparison
has been introduced.

The equalities are valid over every completed residue field of
the original analytic space, not just at rigid points.  They
retain nonproduct base units, moving punctures, all real intercept
walls and arbitrary passive directions.  They impose no extra
growth condition on Green forms.  The unresolved full three-pivot
closed triple remains unresolved: (UTG14) compares two lifts of
one fixed two-pivot program, not the three different programs
whose defect is the finite NFG triple.

\paragraph{Exact dependencies.}
TPG: \path{7fc1508b7ecbcb7ae21ae8cf2be4f5f8712b3bec77e12a43018847e8e87c611e}.
EFT: \path{debbd302e4766854f7512fa75b515571dcf28095434b01120021bd56446ff318}.
EFC: \path{36724216ce00200e1661737a6e105686a7dd90cf80249e8bf9e8bbc516a9dd65}.
XTF: \path{5c6a4b0a99290b0c7dbf9193dfa964c4bcd7988c40fc1376e84b354c28eeebf8},
only its exact hereditary product formula and zero-tail proof.
GNT: \path{c93884ab4c16c71fdf48ff8b4bbae75e71e78c8bf9067dd06bd0a8e0c7b88c25},
only its specified finite full norm/cap program and zero-extension audit.
PTH: \path{b83583762ed294a4365b16dd26272daae1799bad5f59445336a8395507f88652}.
GIC: \path{2ea4db21676b04cc4f2bbebfe168fec80c405f228e546f16550b9445d707c056}.
All native frozen-call hashes and their exact review boundaries are
those of these sources.  No historical file is altered by the
coefficient-uniform strengthening proved here.
